\chapter[2015 --- Prizes]{2015: Multiple Prizes}
\label{ch:2015_Campbell}

\section*{Main Contribution}

\textit{Discovered a novel therapy against parasitic roundworm infections using avermectin, saving millions from river blindness and lymphatic filariasis in developing countries.}

\section{Official Citation}

for their discoveries concerning a novel therapy against infections caused by roundworm parasites

\section{William C. Campbell}

\subsection{Background and Historical Context}

William C. Campbell, an Irish-born parasitologist born in 1930 in Ramelton, County Donegal, Ireland, shared half of the 2015 Nobel Prize in Physiology or Medicine with Satoshi \=Omura ``for their discoveries concerning a novel therapy against infections caused by roundworm parasites.'' Campbell's work, conducted at the Merck Institute for Therapeutic Research in New Jersey, led to the development of ivermectin, a drug that has transformed the treatment of parasitic roundworm infections and has had a significant impact on global public health.

Parasitic roundworm infections (helminthiases) are among the most prevalent and neglected tropical diseases, affecting hundreds of millions of people worldwide, predominantly in tropical and subtropical regions with limited access to clean water, sanitation, and healthcare. The major roundworm infections include onchocerciasis (river blindness), lymphatic filariasis (elephantiasis), strongyloidiasis, and ascariasis. These infections cause chronic morbidity, disability, and death, and have been a major cause of blindness and disfigurement in endemic communities.

Onchocerciasis, caused by the filarial worm \textit{Onchocerca volvulus} and transmitted by blackflies of the genus \textit{Simulium}, was particularly devastating. The disease causes severe itching, skin disfigurement, and blindness, and was responsible for widespread blindness in communities along rivers in West Africa where blackflies bred. Entire villages were abandoned because of the threat of river blindness, and the social and economic impact of the disease was enormous.

Before the discovery of ivermectin, the treatment options for onchocerciasis and other roundworm infections were limited and often dangerous. The only drug available for treating onchocerciasis was diethylcarbamazine (DEC), which had significant side effects and was not effective against all stages of the parasite. Suramin, another antiparasitic drug, was toxic and required intravenous administration. The need for safe and effective treatments for roundworm infections was one of the most pressing unmet needs in global health.

\subsection{The Discovery/Contribution}

The story of ivermectin's development began with the discovery of a novel class of antiparasitic compounds produced by soil bacteria. In the late 1970s, researchers at the Kitasato Institute in Tokyo, Japan, isolated a new strain of soil bacteria, \textit{Streptomyces avermitilis}, from a soil sample collected near a golf course in Ito, Shizuoka Prefecture, Japan. Satoshi \=Omura, a microbiologist at the Kitasato Institute, cultured and characterized this bacterium and provided it to researchers at Merck, including William Campbell, for further study.

Campbell and his colleagues at Merck tested the fermentation products of \textit{Streptomyces avermitilis} against a panel of parasitic organisms and discovered that a group of compounds they named ``avermectins'' had notable antiparasitic activity. Avermectins were highly effective against a range of parasitic roundworms, including the agents causing onchocerciasis, lymphatic filariasis, and strongyloidiasis, as well as ectoparasites such as mites and lice.

Campbell led the effort to develop avermectins into a safe and effective drug for human use. His team modified the chemical structure of the most active avermectin component (avermectin B1) through a process of chemical reduction, producing a derivative they named ``ivermectin.'' Ivermectin had improved antiparasitic activity and a better safety profile than the parent compound.

The mechanism of action of ivermectin involves the potentiation of glutamate-gated chloride channels, which are found in invertebrate nerve and muscle cells but not in vertebrate cells. Ivermectin binds to these channels and causes them to open, allowing chloride ions to flow into the cells, which leads to hyperpolarization and paralysis of the parasite's muscles. This paralysis is irreversible and results in the death of the parasite. Because glutamate-gated chloride channels are absent in vertebrates, ivermectin has a high degree of selectivity for parasites and minimal toxicity to humans.

Clinical trials conducted in the late 1970s and early 1980s demonstrated that ivermectin was highly effective against onchocerciasis, with a single oral dose killing the microfilariae (larval stage) of \textit{Onchocerca volvulus} for up to 12 months. The drug was also effective against a range of other roundworm infections, including strongyloidiasis, ascariasis, and lymphatic filariasis. Ivermectin was approved for human use in 1987 under the brand name Mectizan.

The impact of ivermectin on the treatment of onchocerciasis was transformative. Before ivermectin, there was no safe and effective treatment for the disease, and millions of people in West Africa were at risk of blindness from onchocerciasis. The availability of ivermectin, combined with mass drug administration programs, dramatically reduced the burden of onchocerciasis in endemic communities and enabled the resettlement of villages that had been abandoned due to the disease.

\section{Satoshi \=Omura}

\subsection{Background and Historical Context}

Satoshi \=Omura, a Japanese microbiologist born in 1935 in the Wakayama Prefecture, Japan, shared the 2015 Nobel Prize with William Campbell for their contributions to the development of ivermectin. \=Omura's discovery of \textit{Streptomyces avermitilis}, the soil bacterium that produces avermectins, was the essential first step in the development of one of a major antiparasitic drugs in history.

\=Omura's research career at the Kitasato Institute in Tokyo was devoted to the study of soil microorganisms and their potential for producing therapeutically useful compounds. Soil bacteria, particularly those belonging to the genus \textit{Streptomyces}, are well known for their ability to produce a diverse array of bioactive compounds, including antibiotics, antifungals, and immunosuppressants. \=Omura developed innovative methods for isolating and culturing soil microorganisms, and he systematically screened large numbers of bacterial isolates for their production of bioactive compounds.

\subsection{The Discovery/Contribution}

\=Omura's isolation of \textit{Streptomyces avermitilis} was the result of a systematic and methodical search for novel bioactive compounds produced by soil bacteria. He collected soil samples from a wide range of environments, including forests, farmlands, and urban areas, and used his expertise in microbiology to isolate and culture individual bacterial colonies from these samples. He then tested the fermentation products of each isolate for biological activity against a panel of parasites and other organisms.

The isolation of \textit{Streptomyces avermitilis} from a soil sample near a golf course in Ito, Japan, was a particularly fortunate find. \=Omura recognized the potential of the bioactive compounds produced by this organism and provided the bacterial strain to researchers at Merck for further study and development. This collaboration between Japanese and American scientists exemplified the international nature of drug discovery and the importance of scientific exchange across borders.

\=Omura's contribution to the development of ivermectin was recognized by the Nobel Committee as an essential component of the discovery process. His expertise in microbiology and his systematic approach to the search for novel bioactive compounds enabled the identification of the organism that produced avermectins, which Campbell and his colleagues at Merck subsequently developed into ivermectin.

\subsection{Impact on Modern Medicine}

The impact of ivermectin on global health has been enormous. The drug is estimated to have been administered over 2.5 billion times since its introduction, and it has been the cornerstone of mass drug administration programs for the control and elimination of onchocerciasis and lymphatic filariasis. The Onchocerciasis Control Programme (OCP) in West Africa, launched in 1974, used ivermectin as the primary tool for controlling river blindness, and by the time the program was discontinued in 2002, it had protected an estimated 40 million people from the disease.

The Mectizan Donation Program, established by Merck in 1987, has provided ivermectin free of charge to all countries that need it for the treatment of onchocerciasis and lymphatic filariasis. This program, which is one of the longest-running and most successful drug donation programs in history, has enabled the delivery of billions of treatment courses to hundreds of millions of people in endemic communities.

Ivermectin has also been used effectively against a range of other parasitic infections, including strongyloidiasis, cutaneous larva migrans, and pediculosis (lice infestation). Its activity against ectoparasites has also made it a valuable veterinary medicine, used to treat and prevent parasitic infections in livestock and companion animals.

The discovery of ivermectin and its impact on global health demonstrates the potential of natural product-based drug discovery to address major public health challenges. The identification of avermectins from soil bacteria and their development into a life-saving drug illustrates the importance of basic research in microbiology and the value of international collaboration in drug development.

William C. Campbell and Satoshi \=Omura were awarded the Nobel Prize in Physiology or Medicine in 2015 ``for their discoveries concerning a novel therapy against infections caused by roundworm parasites.'' Their work has saved millions of lives and has transformed the treatment of parasitic diseases worldwide.

\section{Tu Youyou}

\subsection{Background and Historical Context}

Tu Youyou, a Chinese pharmaceutical chemist born in 1930 in Ningbo, Zhejiang Province, China, was awarded the other half of the 2015 Nobel Prize in Physiology or Medicine ``for her discoveries concerning a novel therapy against malaria.'' Her discovery of artemisinin, a potent and fast-acting antimalarial compound derived from the sweet wormwood plant (\textit{Artemisia annua}), has saved millions of lives and is one of a major medical discoveries of the twentieth century.

Malaria is one of the deadliest infectious diseases in human history, having killed more people than any other single infectious agent over the course of human evolution. The disease is caused by protozoan parasites of the genus \textit{Plasmodium}, transmitted to humans through the bites of infected female \textit{Anopheles} mosquitoes. The most deadly species, \textit{Plasmodium falciparum}, causes severe and often fatal malaria, particularly in sub-Saharan Africa, where the disease kills hundreds of thousands of people each year, predominantly young children.

By the mid-twentieth century, the fight against malaria had made significant progress through the use of insecticides to control mosquito populations and chloroquine, a synthetic antimalarial drug that was effective against \textit{P. falciparum}. However, by the 1960s and 1970s, chloroquine-resistant strains of \textit{P. falciparum} had emerged and were spreading throughout Southeast Asia and South America, threatening to undo the gains made in malaria control. The need for new antimalarial drugs that could overcome chloroquine resistance was a matter of urgent public health importance.

\subsection{The Discovery/Contribution}

Tu Youyou's discovery of artemisinin was the result of a secret military research program initiated by the Chinese government during the Vietnam War. In the late 1960s, the Chinese government launched Project 523, a classified research program aimed at finding new treatments for malaria to support North Vietnamese troops who were suffering from the disease. Tu Youyou, a researcher at the China Academy of Chinese Medical Sciences in Beijing, was appointed to lead a group searching for new antimalarial compounds from traditional Chinese herbal medicines.

Tu Youyou and her team systematically reviewed the Chinese medical literature for references to the treatment of malaria-like symptoms. They identified approximately 2,000 traditional Chinese herbal remedies and tested them for antimalarial activity in animal models. Among the plants they tested was sweet wormwood (\textit{Artemisia annua}), which had been used in traditional Chinese medicine for centuries to treat intermittent fevers---a symptom consistent with malaria.

The initial results with sweet wormwood extracts were inconsistent, leading Tu Youyou to re-examine the ancient Chinese medical texts more carefully. She found a reference in a fourth-century text by Ge Hong, a Taoist alchemist, that described the preparation of sweet wormwood for the treatment of intermittent fevers. Crucially, the text specified that the plant material should be soaked in cold water and the extract drunk as a liquid, rather than being boiled or heated as was common in traditional Chinese medicine preparation. This detail suggested that the active compound in sweet wormwood might be heat-sensitive and could be destroyed by boiling.

Following this insight, Tu Youyou modified her extraction method to use low-temperature extraction with diethyl ether, a technique that preserved the activity of the heat-sensitive compound. This modified extraction procedure produced a highly active antimalarial compound, which she named ``qinghaosu'' (artemisinin in English), after the Chinese name for sweet wormwood (qinghao).

Artemisinin had remarkable antimalarial properties. It was rapidly acting, killing \textit{Plasmodium} parasites faster than any existing antimalarial drug. It was effective against chloroquine-resistant strains of \textit{P. falciparum}, and it had a unique mechanism of action involving the formation of free radicals that damaged the parasite's cellular structures.

Tu Youyou personally tested the safety of the compound by volunteering to take the first dose herself, a decision that reflected her commitment to the research and her confidence in the compound's safety. Subsequent clinical trials in malaria patients in China and Vietnam demonstrated that artemisinin was highly effective in curing malaria, with cure rates exceeding 90\%.

\subsection{Impact on Modern Medicine}

The impact of Tu Youyou's discovery of artemisinin on global health has been transformative. Artemisinin-based combination therapies (ACTs) have become the standard treatment for uncomplicated \textit{P. falciparum} malaria worldwide, and are recommended by the World Health Organization (WHO) as the first-line treatment for the disease. ACTs combine artemisinin with a longer-acting partner drug to ensure complete parasite clearance and reduce the risk of drug resistance.

Since their introduction, ACTs have contributed to a dramatic decline in malaria mortality and morbidity. Between 2000 and 2015, the global malaria death rate fell by 60\%, with an estimated 6.2 million lives saved, largely attributable to the widespread use of ACTs. The WHO estimates that artemisinin and its derivatives are responsible for saving hundreds of thousands of lives each year, predominantly in sub-Saharan Africa, where the burden of malaria is notable.

The discovery of artemisinin has also had a significant impact on drug discovery and development. Artemisinin demonstrated the value of natural product-based drug discovery and the potential of traditional medical knowledge to provide leads for the development of modern pharmaceuticals. The success of artemisinin has inspired renewed interest in the study of traditional herbal medicines as sources of novel bioactive compounds.

However, the emergence of artemisinin-resistant strains of \textit{P. falciparum} in Southeast Asia has raised concerns about the long-term effectiveness of ACTs. Artemisinin resistance, characterized by delayed parasite clearance following treatment, has been detected in several countries in the Greater Mekong Subregion and has the potential to spread to Africa, where the malaria burden is notable. The development of new antimalarial drugs and strategies to combat artemisinin resistance is currently a major priority in malaria research.

Tu Youyou was awarded the Nobel Prize in Physiology or Medicine in 2015 ``for her discoveries concerning a novel therapy against malaria.'' She was the first Chinese woman to win a Nobel Prize, and her discovery of artemisinin is one of a major contributions to global health in the history of medicine.

\subsection{Impact on Modern Medicine}

The combined contributions of Campbell, \=Omura, and Tu Youyou have had a transformative impact on the treatment of parasitic diseases worldwide. Ivermectin and artemisinin are two of a major drugs in the history of tropical medicine, and their discovery has saved millions of lives and alleviated untold suffering in some of the world's most vulnerable populations.

Their work demonstrates the power of natural product-based drug discovery and the importance of basic research in addressing global health challenges. The identification of avermectins from soil bacteria and artemisinin from sweet wormwood illustrates the potential of natural products to provide leads for the development of life-saving pharmaceuticals, and highlights the value of both traditional knowledge and modern scientific methods in drug discovery.

The Nobel Committee's recognition of these three laureates in 2015 reflected the global importance of their contributions to medicine and public health. Their discoveries have had a lasting impact on the treatment of parasitic diseases and have contributed to significant reductions in global morbidity and mortality from these conditions.

\section{Legacy}

The legacy of Campbell, \=Omura, and Tu Youyou's work extends across multiple areas of medicine and public health. Their discoveries have provided life-saving treatments for some of the world's most devastating parasitic diseases and have demonstrated the potential of natural product-based drug discovery to address major global health challenges.

William Campbell's work at Merck in developing ivermectin from avermectins represents one of the great success stories of pharmaceutical research. The Mectizan Donation Program, which has provided ivermectin free of charge to endemic communities for over three decades, is a model for drug donation programs and public-private partnerships in global health.

Satoshi \=Omura's discovery of \textit{Streptomyces avermitilis} and his systematic approach to the search for novel bioactive compounds from soil bacteria have inspired generations of microbiologists and natural product chemists. His work has contributed to the broader understanding of the diversity and pharmacological potential of soil microorganisms.

Tu Youyou's discovery of artemisinin is one of a major medical discoveries of the twentieth century. Her use of traditional Chinese medical texts as a source of leads for drug discovery has highlighted the value of traditional knowledge and has inspired renewed interest in the study of herbal medicines. Her story also illustrates the importance of perseverance and innovation in the face of scientific challenges and resource constraints.

Together, the work of these three laureates has saved millions of lives and has transformed the treatment of parasitic diseases worldwide. Their legacy continues to inspire new research and clinical innovation, and their contributions will remain central to the field of tropical medicine for generations to come.


\section{Modern Developments}

Campbell, Ōmura, and Tu Youyou's antiparasitic work has led to modern tropical medicine. Artemisinin-based combination therapies (ACTs) remain the most effective malaria treatment, but artemisinin resistance is emerging in Southeast Asia. New antimalarials (KAF156, cipargamin) are in clinical trials. Understanding of avermectin biology has informed development of new antiparasitic drugs. Understanding of drug resistance mechanisms has led to combination therapy strategies.
